\documentclass[11pt]{article}
\usepackage{fontspec}
\setmainfont{Latin Modern Roman}
\usepackage{amsmath,amssymb,mathtools,unicode-math}
\setmathfont{Latin Modern Math}
\usepackage{xurl,hyperref}
\hypersetup{colorlinks=true,urlcolor=blue,pdftitle={Problem statement}}
\newcommand{\sref}[2]{\hyperlink{source:#1:#2}{[S#2]}}
\newcommand{\cataloguescope}[1]{\paragraph{Mathematical question.}#1}
\begin{document}
\section{Research Problems in Function Theory — Problem 2.90}\label{statement}

\subsection{Definitions and mathematical statement}

\paragraph{Shared notation.}$\mathbb N=\{1,2,\ldots\}$, and $\mathbb Z,\mathbb Q,\mathbb R,\mathbb C$ denote the integers, rationals, reals and complex numbers. Any other conventions are specified locally.


The circle is $\mathbb R/\mathbb Z$. For an orientation-preserving circle homeomorphism $f$ with lift $F(x+1)=F(x)+1$, its rotation number is $\lim_{n\to\infty}(F^{\circ n}(x)-x)/n$ modulo $1$. An irrational $\alpha$ is non-Diophantine if no $C,\tau>0$ satisfy $|\alpha-p/q|\geq Cq^{-(2+\tau)}$ for every $p\in\mathbb Z$, $q\geq1$. An analytic conjugacy is a real-analytic circle diffeomorphism $h$ such that $h\circ f\circ h^{-1}(x)=x+\alpha\pmod1$.
\paragraph{Statement recorded in the supplied source.}
Does there exist a non-Diophantine irrational $\alpha$ for which every orientation-preserving real-analytic circle diffeomorphism of rotation number $\alpha$ admits an analytic conjugacy to the rotation by $\alpha$? \sref{statement}{2}
\subsection{Short English statement}
Is there a non-Diophantine rotation number for which every analytic circle diffeomorphism with that number can be analytically linearized?
\subsection{Sources}
\begin{description}
\item[\hypertarget{source:statement:1}{S1}] ULAM.AI and the UnsolvedMath Contributors, UnsolvedMath: A Curated Collection of Open Mathematics Problems, catalogue entry. CC BY 4.0. \url{https://huggingface.co/datasets/ulamai/UnsolvedMath}.
\item[\hypertarget{source:statement:2}{S2}] W. K. Hayman and E. F. Lingham, \emph{Research Problems in Function Theory}, arXiv:1809.07200v2 (2018), Problem 2.90. \url{https://arxiv.org/abs/1809.07200}.
\end{description}
\label{end:statement}
\end{document}
